\scanpage{236}
On trouve donc un homomorphisme
\[
\varphi:L^\bullet\longrightarrow\mathcal H^\bullet
\]
par
\diagram{\begin{tikzcd}[column sep=large]
\cdots\arrow[r,no head]&0\arrow[r,no head]&\mathcal G\arrow[r,"i"]\arrow[d,"\unsure{\varphi^{-1}}"']&\overline{\mathcal G}\arrow[r,no head]\arrow[d,"\varphi^0"]&0\arrow[r,no head]&\cdots\\
&\mathcal H^{-2}\arrow[r]&\mathcal H^{-1}\arrow[r,"\delta^{-1}",no head]&\mathcal H^0\arrow[r]&\mathcal H^1\arrow[r,no head]&\cdots
\end{tikzcd}}
défini par les conditions
\[
\varphi^{-1}=\unsure{h|k}\qquad\varphi^0\vert\mathcal P=\unsure{u|v}
\]
et la commutativité du carré

[\deleted{\unsure{On voit que}} Un hom de $\overline{\mathcal G}$ dans un un faisceau $F$ $=\{$Hom $v$ de $\mathcal G$ dans $F$, et hom $u$ de $\mathcal P$ dans $F$, avec condition $u(pg)=u(p)v(g)\}$.
Donc prenant $\deleted{\unsure{\delta\varphi^{-1}}}F=\mathcal H^0$, et $u,v$ comme dessus, on trouve $\varphi^0$, avec $\varphi^0i=u$. Mais
\[
u=\delta\deleted{\unsure{\varphi^{-1}}}\unsure{h|k}=\delta\varphi^{-1}\ \ldots
\]
]
Tensorisant avec $\underline O_X$, on trouve
\diagram{\begin{tikzcd}
\underline O_X\arrow[r,"\sim","\substack{\text{dans cat}\\\text{dérivée}}"']&\deleted{\unsure{L^\bullet}}\arrow[r]&\mathcal H\arrow[r]&R\underline{\operatorname{Hom}}(K',K'')
\end{tikzcd}}
dans
\deleted{On a donc}
\[
\boxed{\underline O_X\xrightarrow{\sim}R\underline{\operatorname{Hom}}(K',K'')\quad\text{ds cat. dérivée}}
\]
Donc
\diagram{\begin{tikzcd}
\deleted{\unsure{R\Gamma_X(L^\bullet)}}\arrow[r,"\sim",no head]&R\Gamma_X(\unsure{\underline O_X|\mathbf Z_X})\arrow[r]&\deleted{\unsure{R\Gamma_X\mathcal H}}\arrow[r]&R\operatorname{Hom}(K',K'')
\end{tikzcd}}
d'où
\[
\boxed{R\Gamma_X(\unsure{\underline O_X|\mathbf Z_X})\longrightarrow R\operatorname{Hom}(K',K'')}
\]
Par suite, \deleted{faisant} passant au $R\Gamma_X^0$, on trouve un élément de $R^0\operatorname{Hom}(K',K'')$, i.e. un hom de la cat. dérivée
\margin{235}
